\documentclass[10pt,a4paper]{amsart}
\usepackage[margin=19mm]{geometry}
\usepackage{amsmath,amssymb,mathtools}
\usepackage[scr=rsfs,cal=euler]{mathalfa}
\usepackage{xfrac}
\usepackage[unicode,hidelinks,bookmarks=false]{hyperref}
\newcommand{\ie}{i.e.\ }
\newcommand{\cf}{cf.\ }
\newcommand{\resp}{respectively\ }
\newcommand{\Subsection}{\subsection}
\providecommand{\nobreakdash}{\nobreak\hbox{-}}
\setlength{\emergencystretch}{2em}
\multlinegap0pt
\usepackage{amssymb}
\usepackage{xcolor}
\usepackage[normalem]{ulem}
\newtheorem{lemma}{Lemma}
\title{Additional Congruences for generalized Color Partitions of Hirschhorn and Sellers}
\author{Anjelin Mariya Johnson, James A. Sellers, S.N. Fathima}
\date{Source: arXiv:2510.25250v2 (CC BY 4.0)}
\begin{document}
\maketitle

\noindent\emph{Source and licence.} Additional Congruences for generalized Color Partitions of Hirschhorn and Sellers. Version arXiv:2510.25250v2 (CC BY 4.0), distributed by arXiv under the Creative Commons Attribution 4.0 International licence, http://creativecommons.org/licenses/by/4.0/, as recorded in the arXiv licence metadata for this exact version and shown on the abstract page https://arxiv.org/abs/2510.25250v2. Full licence notice: \texttt{licenses/arxiv-2510.25250-CC-BY-4.0.txt}.\par\smallskip

\noindent\textit{Editorial scope.} Counted unit: the article's lemma l1 (source0.tex lines 248--254). The article's complete original proof of it is delivered with it (source0.tex lines 255--271, one \texttt{proof} environment of 17 lines). No mathematics is written, repaired or re-derived: the statement and the proof are the article's own lines, in the article's own notation, and no retained line is substituted or re-flowed.  No further source-local context is needed: every cross-reference of every retained line resolves inside the two delivered spans.  Nothing was omitted from any retained span, so the package carries no omission record.  No retained line contains a citation, so the wrapper carries no bibliography and no bibliography entry.  No retained line contains a figure environment, an image-inclusion command or drawing code, so no image is delivered and the package declares no asset dependency.  Editorial formatting only, in addition: an AMS \texttt{amsart} preamble that restates the article's own declarations and macros used by the retained lines, namely the environments \texttt{lemma} on separate counters, together with the standard packages those lines need.  The retained environments print in this document as Lemma 1 in order of appearance, whereas the article prints the same items with the numbering of its own full document.  The complete native source of the article as distributed in the arXiv e-print for this exact version is bundled unchanged as \texttt{sources/188017/source0.tex}.
	\begin{lemma}\label{l1}
	We have 
	\begin{align*}
		f_1^3 \equiv \beta_0+\beta_1+\beta_3+\beta_6\pmod{2}
	\end{align*}
where $\beta_i$ consist of terms in which the exponent on $q\equiv i\pmod{7}$.
\end{lemma}

\begin{proof}
	We start by noting that $\frac{k^2+k}{2}\equiv0,1,3 \; and \;\; 6 \; \pmod{7}$, and ${2k+1} \equiv 1 \pmod{2}$.
	Using Jacobi's Theorem, we see that
	\begin{align*}
		(q;q)_\infty^{3} 
		=& \sum_{k=0}^{\infty} (-1)^{k} (2k + 1) q^{\frac{k^{2} + k}{2}} \\
		=& 1 - 3q + 5q^{3} - 7q^{6} + 9q^{10} - 11q^{15} + 13q^{21} - 15q^{28} + 17q^{36}
		- 19q^{45} + - \cdots\\
		\equiv& 1 - q + q^{6} - q^{10} - q^{15} + q^{21} - q^{36} + q^{45} + q^{55} - q^{66} + \cdots \pmod{2}\\
		\equiv &(1 - q^{21} - q^{28} + \cdots)
		- q(1 - q^{14} - q^{35} + \cdots)
		+ q^{3}(1 - q^{7} - q^{42} + \cdots)\\
		&-q^6(1-q^{30}-) \pmod{2}\\
		\equiv& \beta_0+\beta_1+\beta_3+\beta_6\pmod{2}.
	\end{align*}
This completes the proof of Lemma \ref{l1}.
\end{proof}

\end{document}
